\documentclass[11pt]{article}
\usepackage[margin=2.4cm]{geometry}
\usepackage{fontspec}
\setmainfont{DejaVu Sans}
\usepackage{juliamono}
\usepackage{fvextra}
\usepackage{xcolor}
\usepackage{xurl}
\usepackage{hyperref}
\usepackage{amssymb}
\newcommand{\checkbox}{$\square$}
\newcommand{\answerchoices}{\checkbox\ YES \quad \checkbox\ NO \quad \checkbox\ CANNOT DETERMINE}
\newcommand{\decisionexplanation}{%
  \par\smallskip\textbf{If NO, give the reason for the mismatch (required). If
  CANNOT DETERMINE, state what context or evidence is missing:}\par
  \rule{\linewidth}{0.4pt}\par\vspace{7pt}\rule{\linewidth}{0.4pt}\par}
\setlength{\parindent}{0pt}\setlength{\parskip}{5pt}
\setlength{\emergencystretch}{2em}
\DefineVerbatimEnvironment{Vsrc}{Verbatim}{fontsize=\small,frame=single,
  breaklines=true,breakanywhere=true}
\DefineVerbatimEnvironment{Vlean}{Verbatim}{fontsize=\scriptsize,frame=single,
  numbers=left,numbersep=5pt,breaklines=true,breakanywhere=true}
\hypersetup{colorlinks=true,urlcolor=blue}

\begin{document}
\begin{center}{\LARGE Proof Verification Document}\\[2pt]{\large reddit\_askmath\_1kehuzl}\\[2pt]{\small Reddit proof-check posts (newly confirmed)}\end{center}
\begin{center}\fcolorbox{blue!60!black}{blue!4}{\parbox{0.92\linewidth}{\textbf{Make three independent decisions.} Follow the evidence chain: original source $\longrightarrow$ challenged claim $\longrightarrow$ proposed counterexample $\longrightarrow$ Lean theorem. A record is valid only when all three links are correct. Complete the human sign-off before reading the automated review appendix; those earlier results are background, not an answer key.}}\end{center}
\textbf{Use only this document, except for the public source link.} Everything needed for Questions 2 and 3 is included here. For Question 1, open the public source link when one is shown. If there is no public link, the complete stored problem and solution are included here. If the information is still not enough, choose CANNOT DETERMINE.
For each question, choose \textbf{YES}, \textbf{NO}, or \textbf{CANNOT DETERMINE}. Choose CANNOT DETERMINE if the document does not provide enough information for you to decide; do not mark YES if you are uncertain. If you find a specific mismatch, choose NO and explain it in the space below that question. Do not rewrite the claim, add assumptions, or substitute a different counterexample.
\textbf{How copied text is shown.} The problem, quoted passage, and proposed counterexample come from the stored record. The Lean theorem statement and Lean files are copied exactly from the certificate. Some long lines continue on the next printed line; this is only page formatting and does not change the copied text. Line numbers are added only to help you find passages. If the stored record contains a character that cannot be printed, the document shows \texttt{[U+XXXX CONTROL]} in its place.
\begin{description}
\item[Challenged claim.] The exact statement from the source that this record says is false.
\item[Proposed counterexample.] A concrete case that should satisfy the claim's assumptions but make its conclusion false.
\item[Lean certificate.] A Lean file proving a formal theorem. Successful compilation proves the theorem as written, not that it is the right translation of the source.
\end{description}
\section*{Question 1 --- Does this proof verification document accurately represent the source?}
\textbf{What to decide.} Compare the challenged passage below with the original source---either the public link or the complete source text in this document. Answer YES only if they make the same claim and include the same assumptions or conditions. If this document changes the meaning, leaves out a condition, or points to a different passage, answer NO.
\subsection*{Original source}
\textbf{Source record:}
\begin{Vsrc}
qedbench_certificate_availability_triage/data/records_reddit_exhaustive.jsonl
\end{Vsrc}
\textbf{Public source:}\par
\url{https://www.reddit.com/r/askmath/comments/1kehuzl/is_there_a_formalism_where_the_set_of_naturals_is/}\par
\subsection*{Problem or question from the source (verbatim)}
\begin{Vsrc}
Is there a formalism where the set of naturals is twice as big as the set of even naturals?
\end{Vsrc}
\subsection*{Challenged passage recorded from the source (stored text)}
\begin{Vsrc}
between the set of even naturals and the set of sets of two natural numbers.  (n maps to {n, n+1}.)
\end{Vsrc}
\subsection*{Surrounding source lines (stored text with added line labels)}
The labels refer to lines in the solution text stored with the source record; the source-line content is not transliterated.
\begin{Vsrc}
0001 | Context: I studied CS (with the corresponding limited amount of math) ages ago, and I sometimes think about math under the shower...
0002 | 
0003 | I'm not sure if the flair is correct.
0004 | 
0005 | I understand the bijection argument that the set of even naturals has the same cardinality as the set of all naturals.  But it's also just so intuitive to see that the set of even naturals must be half the size as the set of all naturals; after all, every other number isn't even.
0006 | 
0007 | So I tried to come up with some bijections, e.g. between the set of even naturals and the set of sets of two natural numbers.  (n maps to {n, n+1}.)  So since on the right side we always have two natural numbers, on the left hand side we have only one.  But then I thought that it's probably possible to use this to show that there are twice as many naturals as naturals, which doesn't make sense.
0008 | 
0009 | And then it occured to me: for any n, there are n naturals but only n/2 (give or take one) even naturals less than or equal to n.  But it's not clear to me whether this somehow generalizes to a statement about all naturals.  It seems like it should, similar to proof by induction.
0010 | 
0011 | Is there some formalism where the intuitive idea that there are half as many even naturals as naturals?  And are there other interesting results from this formalism?
0012 | 
0013 | I'm happy with a pointer to the right Wikipedia page.  I don't quite know what to search for, though.
\end{Vsrc}
\textbf{A YES means:} the proof verification document faithfully represents both the source context and the challenged claim, without dropping a condition or changing its meaning.
\section*{Question 2 --- Does the proposed counterexample disprove that claim?}
\textbf{Claim to test (repeated for convenience).}
\begin{Vsrc}
between the set of even naturals and the set of sets of two natural numbers.  (n maps to {n, n+1}.)
\end{Vsrc}
\subsection*{How the record says the argument uses this claim}
This is the record's description, not an established fact; check it against the source.
\begin{Vsrc}
The map from the even natural numbers to all two-element sets of natural numbers, defined by n ↦ {n, n+1}, is a bijection.
\end{Vsrc}
\subsection*{Proposed counterexample (exact values or construction)}
\begin{Vsrc}
Let T={0,2}, a two-element set of natural numbers. If T={n,n+1} for an even natural n, then its elements would differ by 1. But the elements of T differ by 2. Equivalently, 0∈T forces n=0 or n+1=0; either case cannot yield {0,2}. Thus T has no preimage, so the map is not surjective and hence not bijective.
\end{Vsrc}
\subsection*{Why the record says the counterexample works}
Check this explanation independently; it is evidence to assess, not an answer key.
\begin{Vsrc}
The codomain contains nonconsecutive pairs, while every image is a pair of consecutive naturals.
\end{Vsrc}
\subsection*{How to decide}
\begin{enumerate}
\item List every assumption that applies to the challenged claim.
\item Check that the proposed example satisfies every one of those assumptions.
\item Recompute the claimed outcome and the actual outcome using exact reasoning.
\item Confirm that the conflict is with the challenged claim itself, not an unrelated sentence.
\end{enumerate}
\textbf{A YES means:} the example is fully specified, satisfies all relevant assumptions, and makes the challenged claim false. If a necessary value, condition, or calculation is missing, choose CANNOT DETERMINE.
\section*{Question 3 --- Does the final Lean theorem statement verify this same counterexample?}
\textbf{What to decide.} A Lean file can compile while proving the wrong theorem. Compare these three links independently:
\begin{enumerate}
\item The Lean definitions represent the same mathematical objects, assumptions, and operations as the challenged claim.
\item The Lean construction uses the same concrete values or example shown in Question 2.
\item The final theorem statement asserts that the challenged claim fails on that example, not merely a weaker or nearby statement.
\end{enumerate}
\subsection*{Selected final Lean theorem statement (verbatim UTF-8 excerpt)}
This is an exact contiguous slice of the certificate file, from the theorem keyword through the original \texttt{:=} token where its proof begins. It is shown first so the formal target is not buried in the complete file.
\begin{Vlean}
theorem refutes_claimed_bijection :
    TwoElem (Pair 0 2) ∧
    (¬ ∃ n : Nat, Even n ∧ SetEq (Image n) (Pair 0 2)) ∧
    ¬ Surjective ∧
    ¬ Bijection :=
\end{Vlean}
\subsection*{Project-local imported Lean files (complete, verbatim UTF-8)}
This section follows the certificate's import commands recursively and prints every repository-owned Lean source file on which the final theorem statement depends. Each file is copied in full, including its definitions and controls. External Mathlib, Std, and Lean Init source is not duplicated.
No project-local imported Lean files were found. External Mathlib, Std, and Lean Init source is not duplicated in this proof verification document.
\subsection*{Complete counterexample Lean file (verbatim UTF-8)}
The complete Lean file is included below so you can check every definition and supporting result used by the final theorem statement. Use it together with the project-local imported Lean files above.
\textbf{Lean certificate file:}
\begin{Vsrc}
partB_feasibility/reddit_certificates/reddit_askmath_1kehuzl.lean
\end{Vsrc}
\begin{Vlean}
/-!
Record: reddit_askmath_1kehuzl
Public URL:
https://www.reddit.com/r/askmath/comments/1kehuzl/is_there_a_formalism_where_the_set_of_naturals_is/

Exact target scope: the proposed map from the even natural numbers to all
two-element sets of natural numbers, n ↦ {n, n+1}, is a bijection.

Concrete refuting codomain witness: {0,2}.  It is a genuine two-element set,
but every image is a consecutive pair.  More explicitly, if an image contains
0, then 0 = n or 0 = n+1.  The successor case is impossible; the first case
forces n = 0, whose image contains 1 although {0,2} does not.  Thus {0,2} has
no preimage, so the proposed map is not surjective and hence not bijective.
The independently verified negative control is the attained pair {2,3}, with
even preimage 2.

Material modelling judgments: this is the actual unbounded `Nat` model, not a
finite proxy.  A set is its membership predicate `Nat → Prop`; `Pair a b`
is the unordered predicate x = a ∨ x = b; and `SetEq` is pointwise logical
equivalence, avoiding function extensionality.  `TwoElem` requires distinct
witnesses and extensional equality with their pair.  `Even n` literally means
∃ k, n = k+k.  `Image n` is `Pair n (n+1)`.  Surjectivity quantifies over
every predicate satisfying `TwoElem` and requires an even preimage with equal
image.  Injectivity is restricted to even inputs, and `Bijection` contains both
that injectivity statement and the stated surjectivity.  The well-definedness
theorem records that every even input lands in the two-element codomain; the
negative control records the real domain condition, codomain condition, and
image equality rather than merely checking an isolated arithmetic identity.

Route review: `SELF_CONTAINED_LEAN` is appropriate because Lean core can
represent the actual infinite source structure directly with unbounded `Nat`
and predicates `Nat → Prop`.  No finite analogue, imported set library, or
weakened witness proposition is needed.  This file intentionally has zero
imports and uses only explicit kernel-checked proofs.
-/

namespace RedditAskmath1kehuzl

def Pair (a b x : Nat) : Prop :=
  x = a ∨ x = b

def SetEq (P Q : Nat → Prop) : Prop :=
  ∀ x, P x ↔ Q x

def TwoElem (P : Nat → Prop) : Prop :=
  ∃ a b, a ≠ b ∧ SetEq P (Pair a b)

def Even (n : Nat) : Prop :=
  ∃ k, n = k + k

def Image (n : Nat) : Nat → Prop :=
  Pair n (n + 1)

def Surjective : Prop :=
  ∀ P : Nat → Prop, TwoElem P →
    ∃ n : Nat, Even n ∧ SetEq (Image n) P

def Injective : Prop :=
  ∀ m n : Nat, Even m → Even n →
    SetEq (Image m) (Image n) → m = n

def Bijection : Prop :=
  Injective ∧ Surjective

theorem nat_ne_add_one (n : Nat) : n ≠ n + 1 := by
  induction n with
  | zero =>
      intro h
      cases h
  | succ n ih =>
      intro h
      apply ih
      exact Nat.succ.inj h

/-- Every input from the stated even-natural domain genuinely lands in the
stated codomain of two-element sets. -/
theorem image_is_two_element (n : Nat) (_hn : Even n) :
    TwoElem (Image n) := by
  refine ⟨n, n + 1, nat_ne_add_one n, ?_⟩
  intro x
  exact Iff.rfl

theorem witness_is_two_element : TwoElem (Pair 0 2) := by
  refine ⟨0, 2, ?_, ?_⟩
  · intro h
    cases h
  · intro x
    exact Iff.rfl

/- Core does not need a synthesized `Decidable` instance for this custom
predicate: the two impossible numeral equalities are eliminated explicitly. -/
theorem one_is_not_in_witness : ¬ Pair 0 2 1 := by
  intro h
  cases h with
  | inl oneEqZero =>
      cases oneEqZero
  | inr oneEqTwo =>
      cases oneEqTwo

theorem witness_has_no_even_preimage :
    ¬ ∃ n : Nat, Even n ∧ SetEq (Image n) (Pair 0 2) := by
  intro h
  cases h with
  | intro n hPreimage =>
      cases hPreimage with
      | intro hn hEq =>
          have imageContainsZero : Image n 0 :=
            (hEq 0).mpr (Or.inl rfl)
          cases imageContainsZero with
          | inl zeroEqN =>
              cases zeroEqN
              have targetContainsOne : Pair 0 2 1 :=
                (hEq 1).mp (Or.inr rfl)
              exact one_is_not_in_witness targetContainsOne
          | inr zeroEqSucc =>
              cases zeroEqSucc

theorem claimed_surjectivity_fails : ¬ Surjective := by
  intro hSurjective
  exact witness_has_no_even_preimage
    (hSurjective (Pair 0 2) witness_is_two_element)

/-- Load-bearing refutation of the literal published bijection claim: the
codomain witness {0,2} is valid, has no even preimage, and therefore falsifies
both the full surjectivity proposition and the stated bijection proposition. -/
theorem refutes_claimed_bijection :
    TwoElem (Pair 0 2) ∧
    (¬ ∃ n : Nat, Even n ∧ SetEq (Image n) (Pair 0 2)) ∧
    ¬ Surjective ∧
    ¬ Bijection := by
  refine ⟨witness_is_two_element, witness_has_no_even_preimage,
    claimed_surjectivity_fails, ?_⟩
  intro hBijection
  exact claimed_surjectivity_fails hBijection.2

/-- Meaningful negative control: in the same unbounded model, the even input
2 really attains the valid two-element codomain member {2,3}.  This detects a
uniformly false image relation or a vacuous domain/codomain encoding. -/
theorem negative_control_consecutive_pair_is_attained :
    Even 2 ∧
    TwoElem (Pair 2 3) ∧
    SetEq (Image 2) (Pair 2 3) := by
  refine ⟨⟨1, rfl⟩, ?_, ?_⟩
  · refine ⟨2, 3, nat_ne_add_one 2, ?_⟩
    intro x
    exact Iff.rfl
  · intro x
    exact Iff.rfl

end RedditAskmath1kehuzl

\end{Vlean}
\subsection*{Proof-safety information}
The automated scan found no unfinished proofs, new axiom declarations, unsafe declarations, external code hooks, or uses of Lean's compiled decision procedure. This scan concerns proof construction only; it does not establish that the theorem matches the source claim.
\textit{Technical scan counts:} native\_decide=0, sorry=0, axiom\_decl=0, unsafe=0, extern=0. \textit{Technical tier:} KERNEL-ONLY.
\textbf{A YES means:} Lean represents the same objects and assumptions, uses the same proposed counterexample, and the final theorem statement asserts the same contradiction. Successful compilation by itself is not enough.
\clearpage
\section*{Human sign-off}
Complete this block from your own review of Questions 1--3. Do not use the automated appendix below as an answer key.\par
\textbf{Reviewer name:} \rule{9cm}{0.4pt}\par
\textbf{Date:} \rule{4cm}{0.4pt}\par
\medskip\textbf{1. Does the proof verification document accurately represent the source and challenged claim?}\par
\answerchoices\par
\decisionexplanation
\medskip\textbf{2. Does the proposed counterexample satisfy the assumptions and make that claim false?}\par
\answerchoices\par
\decisionexplanation
\medskip\textbf{3. Does the final Lean theorem statement represent and verify that same counterexample?}\par
\answerchoices\par
\decisionexplanation
\medskip\textbf{Overall decision:}\par
\checkbox\ VALIDATED (all three answers are YES)\par
\checkbox\ REJECTED (at least one answer is NO)\par
\checkbox\ NEEDS ANOTHER REVIEWER (at least one CANNOT DETERMINE and no NO)\par
\medskip\textbf{Additional comments.} Use this space for any extra source, counterexample, calculation, or Lean details that support your decisions.\par\rule{\linewidth}{0.4pt}\par\vspace{7pt}\rule{\linewidth}{0.4pt}
\clearpage
\section*{Appendix --- Prior automated review (background only)}
\textbf{These results are not an answer key.} They come from the earlier Opus/Fable verification campaign and are printed only so an auditor can see what was previously checked. The human decisions above must be made independently. Raw status codes are retained in small type for traceability.
\textbf{Overall automated disposition:} Cleared by automated screening; an independent search also found a counterexample.
\par\textit{Audit code:} \nolinkurl{PASS_CORROBORATED}
\subsection*{Independent search with no access to the stored example}
\textbf{Result:} An independent search found a counterexample without seeing the stored example.
\par\textit{Audit code:} \nolinkurl{WITNESS_FOUND}
\par\textbf{Checker's explanation:}
\begin{Vsrc}
{0, 2} (with 0 in N; {1, 2} if N starts at 1). It is a set of two natural numbers, hence an element of the claimed codomain, but it has no preimage under f: f(n) = {n, n+1} = {0, 2} forces n = 0 and n + 1 = 2, i.e. 1 = 2, false. Since min f(n) = n, the ONLY candidate preimage of a pair {a, b} with a < b is n = a, so this is a complete argument, not a bounded-window artifact. Hence f is not surjective and therefore not a bijection. (Shortened here; the complete note remains in the machine-check JSON.)
\end{Vsrc}
\par\textbf{Cases checked:}
\begin{Vsrc}
Deterministic stdlib-only exhaustive sweep, N = 200, both natural-number conventions (N containing 0 and N starting at 1). Domain: every even n in [n_min, 200] (101 resp. 100 values); injectivity checked over all of them. Codomain (wide reading): every 2-element subset {a,b} of {n_min,...,200} with a < b (20100 resp. 19900 pairs), each tested for membership in the image of f. (Shortened here; the complete note remains in the machine-check JSON.)
\end{Vsrc}
\subsection*{Independent re-check of the proposed counterexample}
\textbf{Result:} A separate checker recomputed the proposed counterexample and confirmed it.
\par\textit{Audit code:} \nolinkurl{CONFIRMED}
\par\textbf{Checker's note:}
\begin{Vsrc}
The witness is fully specified and finitary; nothing had to be guessed. One strained alternative reading was considered and rejected as not what the quote says: if 'sets of two natural numbers' were silently narrowed to 'sets of two CONSECUTIVE naturals', {0,2} would not be in the codomain -- but even then the map from EVEN n only would still fail surjectivity (it misses {1,2}), so the bijection claim is false on either reading; the recorded witness {0,2} is the correct one for the quote as written. (Shortened here; the complete note remains in the machine-check JSON.)
\end{Vsrc}
\subsection*{Independent review of the source-to-Lean correspondence}
\textbf{Result:} A separate reviewer judged that Lean proves the same failure using the same example.
\par\textit{Audit code:} \nolinkurl{MATCHES}
\par\textbf{Theorem summary from that reviewer:}
\begin{Vsrc}
`refutes_claimed_bijection` is a four-way conjunction over the unbounded-Nat model with sets encoded as membership predicates `Nat -> Prop`: (i) `TwoElem (Pair 0 2)` -- the witness {0,2} really is a two-element set of naturals (distinct witnesses 0!= 2, pointwise-equivalent to the pair predicate); (ii) `NOT EXISTS n, Even n AND SetEq (Image n) (Pair 0 2)` -- no even natural has image {n, n+1} equal to {0,2}; (Shortened here; the complete note remains in the machine-check JSON.)
\end{Vsrc}
\par\textbf{Difference noted by that reviewer:}
\begin{Vsrc}
None material. Link 1 (definitions): `Even n := EXISTS k, n = k + k` is literal evenness; `Pair a b x := x = a OR x = b` with `SetEq` as pointwise iff encodes unordered sets without invoking funext or Classical; `TwoElem` demands two distinct elements, so the codomain is exactly the two-element subsets of Nat, neither weakened (no multisets/ordered pairs) nor strengthened (no bound on the naturals -- this is the genuine infinite Nat, not a finite proxy); (Shortened here; the complete note remains in the machine-check JSON.)
\end{Vsrc}
\subsection*{Follow-up review}
No follow-up review was needed or recorded.
\textit{This appendix is only background. You do not need any separate machine-output files to answer the three questions.}
\end{document}